% GENERATED FILE. Edit canonical lessons, then run npm run generate:books.
% canonical-source-hash: fnv1a64:d795a1435760ef98
% canonical-lessons: RU-C432-protekat, RU-C432-kapat, RU-C432-zasoryatsya, RU-C432-tech, RU-C432-potop

\chapter{A leaking tap}
\label{ch:ru-a-leaking-tap}
\hlchaptermodality{432}

\begin{chapteropening}
\textbf{By the end of this chapter:} \emph{I can report water problems: a leak, dripping, a blocked drain, water running and a flood.}

Complete the last lesson of chapter 432: I can report water problems: a leak, dripping, a blocked drain, water running and a flood.
\end{chapteropening}

\section[\texorpdfstring{\ru{протекать} (protekát)}{протекать (protekát)}]{\ru{протекать} (protekát) — to leak (of a roof or a tap); to flow through}

\label{lesson:RU-C432-protekat}



\begin{quote}
Before the new one: say the Russian for to air, then the Russian for to dry.
\end{quote}

\subsection*{You'll want to know: \ru{протекать}}
\textbf{\ru{протекать}} (\emph{protekát}) — \textquotedblleft{}to leak (of a roof or a tap); to flow through\textquotedblright{}.

\begin{quote}
\ru{В} \ru{ванной} \ru{протекает} \ru{кран}.
\end{quote}

\begin{quote}
\emph{V vánnoy protekáyet kran.}
\end{quote}

\begin{quote}
The tap in the bathroom is leaking.
\end{quote}

Imperfective; perfective \emph{\ru{протечь}}. \emph{\ru{Крыша} \ru{протекает}} — the roof leaks.

\begin{cousinweb}
\emph{\ru{протекать}} is built from \emph{\ru{про}-} 'through' and \emph{-\ru{текать}}, from \emph{\ru{течь}} 'to flow' (EtymDB lists \emph{\ru{течь}} with that sense).
\end{cousinweb}

\begin{grammarlens}[title={What you've built}]
One more word for telling, explaining and describing at B1.
\end{grammarlens}

\subsection*{Your turn}
\begin{itemize}
\raggedright
  \item \emph{Say it:} \emph{protekát}
  \item \emph{Say it:} \emph{protekát}, once more
  \item \emph{Say it:} a sentence of your own with \emph{protekát}
  \item \emph{Recall:} say the Russian for to air, then the Russian for to dry, then say \emph{protekát} again
\end{itemize}

\begin{tcolorbox}[breakable,colback=teal!4,colframe=teal!35!black,title={Before you move on}]
What does \ru{протекать} mean? (\textquotedblleft{}To leak (of a roof or a tap); to flow through\textquotedblright{}.) Say it once more.
\end{tcolorbox}
\section[\texorpdfstring{\ru{капать} (kápat)}{капать (kápat)}]{\ru{капать} (kápat) — to drip}

\label{lesson:RU-C432-kapat}



\begin{quote}
Before the new one: say the Russian for to dry, then the Russian for to leak.
\end{quote}

\subsection*{You'll want to know: \ru{капать}}
\textbf{\ru{капать}} (\emph{kápat}) — \textquotedblleft{}to drip\textquotedblright{}.

\begin{quote}
\ru{Всю} \ru{ночь} \ru{из} \ru{крана} \ru{капала} \ru{вода}.
\end{quote}

\begin{quote}
\emph{Vsyu noch iz krána kápala vodá.}
\end{quote}

\begin{quote}
All night water dripped from the tap.
\end{quote}

Imperfective; perfective \emph{\ru{капнуть}}. Present \emph{\ru{капает}}.

\begin{cousinweb}
\emph{\ru{капать}} goes with \emph{\ru{капля}} 'a drop'; EtymDB lists \emph{\ru{капать}} without a deeper source.
\end{cousinweb}

\begin{grammarlens}[title={What you've built}]
One more word for telling, explaining and describing at B1.
\end{grammarlens}

\subsection*{Your turn}
\begin{itemize}
\raggedright
  \item \emph{Say it:} \emph{kápat}
  \item \emph{Say it:} \emph{kápat}, once more
  \item \emph{Say it:} a sentence of your own with \emph{kápat}
  \item \emph{Recall:} say the Russian for to dry, then the Russian for to leak, then say \emph{kápat} again
\end{itemize}

\begin{tcolorbox}[breakable,colback=teal!4,colframe=teal!35!black,title={Before you move on}]
What does \ru{капать} mean? (\textquotedblleft{}To drip\textquotedblright{}.) Say it once more.
\end{tcolorbox}
\section[\texorpdfstring{\ru{засоряться} (zasoryátsya)}{засоряться (zasoryátsya)}]{\ru{засоряться} (zasoryátsya) — to get blocked, to clog}

\label{lesson:RU-C432-zasoryatsya}



\begin{quote}
Before the new one: say the Russian for to leak, then the Russian for to drip.
\end{quote}

\subsection*{You'll want to know: \ru{засоряться}}
\textbf{\ru{засоряться}} (\emph{zasoryátsya}) — \textquotedblleft{}to get blocked, to clog\textquotedblright{}.

\begin{quote}
\ru{В} \ru{старых} \ru{домах} \ru{часто} \ru{засоряются} \ru{трубы}.
\end{quote}

\begin{quote}
\emph{V stárykh domákh chásto zasoryáyutsya trúby.}
\end{quote}

\begin{quote}
In old houses the pipes often get blocked.
\end{quote}

Imperfective; perfective \emph{\ru{засориться}}. \emph{\ru{Раковина} \ru{засорилась}} — the sink is blocked.

\begin{cousinweb}
\emph{\ru{засоряться}} is built from \emph{\ru{за}-}, \emph{\ru{сор}} 'litter, rubbish' and \emph{-\ru{ся}}: to fill up with rubbish. EtymDB lists \emph{\ru{сор}} without a source.
\end{cousinweb}

\begin{grammarlens}[title={What you've built}]
One more word for telling, explaining and describing at B1.
\end{grammarlens}

\subsection*{Your turn}
\begin{itemize}
\raggedright
  \item \emph{Say it:} \emph{zasoryátsya}
  \item \emph{Say it:} \emph{zasoryátsya}, once more
  \item \emph{Say it:} a sentence of your own with \emph{zasoryátsya}
  \item \emph{Recall:} say the Russian for to leak, then the Russian for to drip, then say \emph{zasoryátsya} again
\end{itemize}

\begin{tcolorbox}[breakable,colback=teal!4,colframe=teal!35!black,title={Before you move on}]
What does \ru{засоряться} mean? (\textquotedblleft{}To get blocked, to clog\textquotedblright{}.) Say it once more.
\end{tcolorbox}
\section[\texorpdfstring{\ru{течь} (tech)}{течь (tech)}]{\ru{течь} (tech) — to flow; to leak}

\label{lesson:RU-C432-tech}



\begin{quote}
Before the new one: say the Russian for to drip, then the Russian for to get blocked.
\end{quote}

\subsection*{You'll want to know: \ru{течь}}
\textbf{\ru{течь}} (\emph{tech}) — \textquotedblleft{}to flow; to leak\textquotedblright{}.

\begin{quote}
\ru{Под} \ru{раковиной} \ru{течёт} \ru{вода}.
\end{quote}

\begin{quote}
\emph{Pod rákovinoy techyót vodá.}
\end{quote}

\begin{quote}
Water is leaking under the sink.
\end{quote}

Imperfective. Present \emph{\ru{теку}, \ru{течёшь}, \ru{текут}}; past \emph{\ru{тёк}, \ru{текла}}. As a feminine noun, \emph{\ru{течь}} is a leak.

\begin{cousinweb}
\emph{\ru{течь}} 'to flow' is usually counted an inherited Slavic verb; EtymDB lists it with that sense but without a source. \emph{\ru{Восток}} 'east' is often linked with the same root.
\end{cousinweb}

\begin{grammarlens}[title={What you've built}]
One more word for telling, explaining and describing at B1.
\end{grammarlens}

\subsection*{Your turn}
\begin{itemize}
\raggedright
  \item \emph{Say it:} \emph{tech}
  \item \emph{Say it:} \emph{tech}, once more
  \item \emph{Say it:} a sentence of your own with \emph{tech}
  \item \emph{Recall:} say the Russian for to drip, then the Russian for to get blocked, then say \emph{tech} again
\end{itemize}

\begin{tcolorbox}[breakable,colback=teal!4,colframe=teal!35!black,title={Before you move on}]
What does \ru{течь} mean? (\textquotedblleft{}To flow; to leak\textquotedblright{}.) Say it once more.
\end{tcolorbox}
\section[\texorpdfstring{\ru{потоп} (potóp)}{потоп (potóp)}]{\ru{потоп} (potóp) — a flood (in a flat); a deluge}

\label{lesson:RU-C432-potop}



\begin{quote}
Before the new one: say the Russian for to get blocked, then the Russian for to flow.
\end{quote}

\subsection*{You'll want to know: \ru{потоп}}
\textbf{\ru{потоп}} (\emph{potóp}) — \textquotedblleft{}a flood (in a flat); a deluge\textquotedblright{}.

\begin{quote}
\ru{Соседи} \ru{сверху} \ru{устроили} \ru{нам} \ru{потоп}.
\end{quote}

\begin{quote}
\emph{Sosédi svérkhu ustróili nam potóp.}
\end{quote}

\begin{quote}
The upstairs neighbours flooded our flat.
\end{quote}

Masculine; colloquial for a flooded flat. \emph{\ru{Всемирный} \ru{потоп}} — the Flood.

\begin{cousinweb}
\emph{\ru{потоп}} is built on \emph{\ru{потопить}} 'to flood, to sink', from \emph{\ru{топить}} 'to drown', which EtymDB takes from Proto-Slavic \emph{topiti}.
\end{cousinweb}

\begin{grammarlens}[title={What you've built}]
One more word for telling, explaining and describing at B1.
\end{grammarlens}

\subsection*{Your turn}
\begin{itemize}
\raggedright
  \item \emph{Say it:} \emph{potóp}
  \item \emph{Say it:} \emph{potóp}, once more
  \item \emph{Say it:} a sentence of your own with \emph{potóp}
  \item \emph{Recall:} say the Russian for to get blocked, then the Russian for to flow, then say \emph{potóp} again
\end{itemize}

\begin{tcolorbox}[breakable,colback=teal!4,colframe=teal!35!black,title={Before you move on}]
What does \ru{потоп} mean? (\textquotedblleft{}A flood (in a flat); a deluge\textquotedblright{}.) Say it once more.
\end{tcolorbox}
